je pro \acacc{ČR} důležitý jako~nejbližší průmyslový a~obchodní referenční systém, v~němž odvětvové vyjednávání a~podnikové spolurozhodování dlouhodobě stabilizovaly mzdový vývoj. Zadruhé německá zkušenost s~reformami Hartz ukazuje, že deregulace okrajových forem zaměstnání bez dostatečné tarifní vazby může vytvořit rozsáhlý nízkopříjmový segment s~následnými socioekonomickými dopady i~v~ekonomice s~vysokou produktivitou.~\cite{JdeToPruzne_PS2_NemeckyModel_2019}\par

Součástí erozního mechanismu je i~tzv.~\enquote{OT-členství} (\emph{Ohne Tarifbindung}), tedy členství zaměstnavatele ve svazu bez tarifní vázanosti, které oslabuje dosah odvětvových smluv i~v~institucionálně silném systému.~\cite{Hala_SystemySocialnihoDialogu2013}\par

Přenositelnou částí není jen samotné odvětvové vyjednávání, ale také podniková participace. Česká rada zaměstnanců existuje právně, avšak nemá váhu německého \textit{Betriebsrat}. Posílení informačních a~konzultačních práv při restrukturalizaci, zavádění automatizace a~při změnách organizace práce by mohlo českým podnikům usnadnit zvládání strukturálních změn. Německý příklad zároveň varuje, že flexibilizace práce musí být doprovázena silnějším \acs{KV}, jinak se z~ní stává cesta k~mzdové segmentaci.\par
